% Options for packages loaded elsewhere
\PassOptionsToPackage{unicode}{hyperref}
\PassOptionsToPackage{hyphens}{url}
\PassOptionsToPackage{dvipsnames,svgnames,x11names}{xcolor}
\documentclass[
  10pt,
  fontset=fandol]{ctexart}
\usepackage{xcolor}
\usepackage[margin=16mm]{geometry}
\usepackage{amsmath,amssymb}
\setcounter{secnumdepth}{-\maxdimen} % remove section numbering
\usepackage{iftex}
\ifPDFTeX
  \usepackage[T1]{fontenc}
  \usepackage[utf8]{inputenc}
  \usepackage{textcomp} % provide euro and other symbols
\else % if luatex or xetex
  \usepackage{unicode-math} % this also loads fontspec
  \defaultfontfeatures{Scale=MatchLowercase}
  \defaultfontfeatures[\rmfamily]{Ligatures=TeX,Scale=1}
\fi
\usepackage{lmodern}
\ifPDFTeX\else
  % xetex/luatex font selection
\fi
% Use upquote if available, for straight quotes in verbatim environments
\IfFileExists{upquote.sty}{\usepackage{upquote}}{}
\IfFileExists{microtype.sty}{% use microtype if available
  \usepackage[]{microtype}
  \UseMicrotypeSet[protrusion]{basicmath} % disable protrusion for tt fonts
}{}
\makeatletter
\@ifundefined{KOMAClassName}{% if non-KOMA class
  \IfFileExists{parskip.sty}{%
    \usepackage{parskip}
  }{% else
    \setlength{\parindent}{0pt}
    \setlength{\parskip}{6pt plus 2pt minus 1pt}}
}{% if KOMA class
  \KOMAoptions{parskip=half}}
\makeatother
\setlength{\emergencystretch}{3em} % prevent overfull lines
\providecommand{\tightlist}{%
  \setlength{\itemsep}{0pt}\setlength{\parskip}{0pt}}
\usepackage{amsmath,amssymb,mathtools,needspace}
\xeCJKDeclareCharClass{CJK}{"2032,"0400->"04FF,"2160->"216B,"2460->"2473,"25A0,"25C7,"25CB,"25CF}
\IfFontExistsTF{Arial Unicode MS}{\xeCJKsetup{AutoFallBack=true}\setCJKfallbackfamilyfont{\CJKrmdefault}{Arial Unicode MS}}{}
\setcounter{secnumdepth}{0}
\setlength{\emergencystretch}{2em}
\allowdisplaybreaks
\usepackage{bookmark}
\IfFileExists{xurl.sty}{\usepackage{xurl}}{} % add URL line breaks if available
\urlstyle{same}
\hypersetup{
  pdftitle={习题},
  colorlinks=true,
  linkcolor={Maroon},
  filecolor={Maroon},
  citecolor={Blue},
  urlcolor={Blue},
  pdfcreator={LaTeX via pandoc}}

\title{习题}
\author{}
\date{}

\begin{document}
\maketitle

\subsection{习题}\label{ux4e60ux9898}

\textbf{1.} 用幂法计算下列矩阵的主特征值及对应的特征向量：

\[
\text{（1）}\ A_1=\begin{pmatrix}7&3&-2\\3&4&-1\\-2&-1&3\end{pmatrix};\qquad
\text{（2）}\ A_2=\begin{pmatrix}3&-4&3\\-4&6&3\\3&3&1\end{pmatrix},
\]

当特征值有三位小数稳定时迭代终止。

\textbf{2.} 方阵 \(T\) 分块形式为

\[
T=\begin{pmatrix}
T_{11}&T_{12}&\cdots&T_{1n}\\
&T_{22}&\cdots&T_{2n}\\
&&\ddots&\vdots\\
&&&T_{nn}
\end{pmatrix},
\]

其中 \(T_{ii}\ (i=1,2,\cdots,n)\) 为方阵，\(T\) 称为块上

\href{../../source-images/pdf-260.jpeg}{原页}

三角阵，如果对角块的阶数至多不超过 2，则称 \(T\) 为准三角形形式。用 \(\sigma(T)\) 记矩阵 \(T\) 的特征值集合，证明 \(\sigma(T)=\displaystyle\bigcup_{i=1}^{n}\sigma(T_{ii})\)。

\textbf{3.} 利用反幂法求矩阵 \(\begin{pmatrix}6&2&1\\2&3&1\\1&1&1\end{pmatrix}\) 的最接近于 6 的特征值及对应的特征向量。

\textbf{4.} 求矩阵 \(\begin{pmatrix}4&0&0\\0&3&1\\0&1&3\end{pmatrix}\) 与特征值 4 对应的特征向量。

\textbf{5.} （1）设 \(A\) 是对称矩阵，\(\lambda\) 和 \(x\ (\|x\|_2=1)\) 是 \(A\) 的一个特征值及相应的特征向量，又设 \(P\) 为一个正交阵，使 \(Px=e_1=(1,0,\cdots,0)^{\mathrm T}\)，证明 \(B=PAP^{\mathrm T}\) 的第 1 行和第 1 列除了 \(\lambda\) 之外其余元素均为零。

（2）对于矩阵 \(A=\begin{pmatrix}2&10&2\\10&5&-8\\2&-8&11\end{pmatrix}\)，\(\lambda=9\) 是其特征值，\(x=\left(\dfrac23,\dfrac13,\dfrac23\right)^{\mathrm T}\) 是相应于 9 的特征向量，试求一初等反射阵 \(P\)，使 \(Px=e_1\)，并计算 \(B=PAP^{\mathrm T}\)。

\textbf{6.} 利用初等反射阵将 \(A=\begin{pmatrix}1&3&4\\3&1&2\\4&2&1\end{pmatrix}\) 正交相似约化为对称三对角阵。

\textbf{7.} 设 \(A\in\mathbb R^{n\times n}\)，且 \(a_{i1},a_{j1}\) 不全为零，\(P_{ij}\) 为使 \(a_{j1}^{(2)}=0\) 的平面旋转阵，试推导 \(P_{ij}A\) 第 \(i\) 行、第 \(j\) 行元素的计算公式及 \(AP_{ij}^{\mathrm T}\) 第 \(i\) 列、第 \(j\) 列元素的计算公式。

\textbf{8.} 设 \(A_{n-1}\) 是由 Householder 方法得到的矩阵，又设 \(y\) 是 \(A_{n-1}\) 的一个特征向量。

（1）证明矩阵 \(A\) 对应的特征向量是 \(x=P_1P_2\cdots P_{n-2}y\)；（2）对于给出的 \(y\) 应如何计算 \(x\)？

\textbf{9.} 用带位移的 QR 方法计算下列矩阵的全部特征值：

\[
\text{（1）}\ A=\begin{pmatrix}1&2&0\\2&-1&1\\0&1&3\end{pmatrix};\qquad
\text{（2）}\ B=\begin{pmatrix}3&1&0\\1&2&1\\0&1&1\end{pmatrix}.
\]

\textbf{10.} 试用初等反射阵将

\[
A=\begin{pmatrix}1&1&1\\2&-1&-1\\2&-4&5\end{pmatrix}
\]

分解为 \(QR\)，其中 \(Q\) 为正交阵，\(R\) 为上三角阵。

\href{../../source-images/pdf-261.jpeg}{原页}

\section{下篇　高效算法设计}\label{ux4e0bux7bc7-ux9ad8ux6548ux7b97ux6cd5ux8bbeux8ba1}

\end{document}
